A equação (2.2) do livro-texto descreve o período de um pêndulo sujeito a oscilações de pequeno ângulo e é dada por:
$$ T_0 = 2\pi\sqrt{\frac{l}{g}} $$ [2]

O objetivo é chegar a esta mesma relação funcional utilizando o \textbf{método básico da análise dimensional}, que consiste em construir uma equação funcional genérica e eliminar suas dimensões primárias passo a passo [3, 4].

\vspace{0.3cm}
\textbf{Passo 1: Equação funcional e variáveis} \\
Começamos listando as variáveis e parâmetros e assumindo que o período do pêndulo ($T_0$) é função do seu comprimento ($l$), da aceleração da gravidade ($g$) e, como uma hipótese inicial intuitiva, da massa do corpo do pêndulo ($m$). A equação funcional inicial é:
$$ T_0 = f(l, g, m) $$

\textbf{Passo 2: Identificar as dimensões físicas} \\
As dimensões físicas fundamentais das variáveis envolvidas são:
\begin{itemize}
    \item Período $[T_0] = \text{T}$ (Tempo)
    \item Comprimento $[l] = \text{L}$ (Comprimento)
    \item Aceleração da gravidade $[g] = \text{L/T}^2$
    \item Massa $[m] = \text{M}$ (Massa)
\end{itemize}

\textbf{Passo 3: Eliminar a dimensão de Massa (M)} \\
Observando a lista de dimensões, notamos que a dimensão primária de Massa ($M$) aparece em apenas uma variável em todo o sistema: a própria massa $m$. Semelhante ao que ocorre no exemplo da queda livre do livro-texto, é impossível cancelar essa dimensão usando as outras grandezas disponíveis [5]. Portanto, para que a equação seja dimensionalmente homogênea, a massa simplesmente não pode ser uma variável atuante neste problema [5]. A nossa equação funcional revisada torna-se:
$$ T_0 = f_1(l, g) $$

\textbf{Passo 4: Eliminar a dimensão de Tempo (T)} \\
A dimensão de tempo ($T$) aparece diretamente no período ($T_0$) e na gravidade ($g$) [6]. Para eliminar o tempo, podemos multiplicar a variável $T_0$ pela raiz quadrada de $g$ [6]. 
Analisando as dimensões dessa nova operação:
$$ \left[ T_0\sqrt{g} \right] = \text{T} \cdot \sqrt{\frac{\text{L}}{\text{T}^2}} = \text{T} \cdot \frac{\sqrt{\text{L}}}{\text{T}} = \sqrt{\text{L}} = \text{L}^{1/2} $$
Como o tempo foi perfeitamente eliminado, a nova quantidade intermediária tem dimensão restando apenas em termos de comprimento [6]. A equação revisada é:
$$ T_0\sqrt{g} = f_2(l) $$

\textbf{Passo 5: Eliminar a dimensão de Comprimento (L)} \\
O lado esquerdo da nossa equação ($T_0\sqrt{g}$) agora possui dimensão puramente de $\sqrt{\text{L}}$ (ou $L^{1/2}$) [6]. A variável que nos resta no lado direito é $l$, que tem dimensão $\text{L}$. Para eliminar completamente as dimensões do problema, dividimos a grandeza que construímos no passo anterior pela raiz quadrada do comprimento ($\sqrt{l}$):
$$ \left[ \frac{T_0\sqrt{g}}{\sqrt{l}} \right] = \frac{\sqrt{\text{L}}}{\sqrt{\text{L}}} = 1 $$

\vspace{0.3cm}
\textbf{Conclusão:} \\
Ao realizarmos esses passos, nós eliminamos todas as dimensões físicas e obtivemos um único grupo inteiramente adimensional. A análise dimensional nos ensina que, quando resta apenas um grupo adimensional no sistema, este grupo deve ser numericamente igual a uma constante [6]. Portanto:
$$ T_0 \sqrt{\frac{g}{l}} = \text{constante} \implies T_0 = \text{constante} \times \sqrt{\frac{l}{g}} $$

A aplicação elementar do método básico confirma perfeitamente a estrutura algébrica da equação (2.2). A análise dimensional revela que o período é rigorosamente proporcional a $\sqrt{l/g}$, embora (como de costume) não possa nos fornecer o valor escalar da constante, que a física analítica demonstra ser $2\pi$ [2, 6].